\section{Thermodynamic model and transferred parameters}\label{sec:p2-model}
\subsection{Shared terms and amine-specific chemistry}
% P:p2-model-01
The residual Helmholtz energy retains the hard-chain, dispersion, association, Debye--H\"uckel and Born contributions of Part I \cite[Sect.~2]{PolleyPartI}.
The water and \COtwo component descriptions, hydronium, hydroxide, bicarbonate and carbonate, all interactions among these six shared species, and the completed water/carbonate reaction laws are held fixed.
Table~\ref{tab:p2-origins} distinguishes independent fits, literature values, parameters estimated jointly with MEA, and assumptions.
Free interactions between an amine-specific ion and a shared species remain amine-specific parameters; freezing bicarbonate does not freeze every interaction involving bicarbonate.

% Generator: docs/scientific/latex/scripts/part2_displays.py; no model evaluation.
% Input: analyses/amp_transfer/inputs/kundu/current-declarations.toml SHA256 3085081adcf6b1d5e797a49ae6b96037cc1b5bc5191f45cff699d8546d774c55
% Input: analyses/amp_transfer/inputs/redesign/refit-306-floor-coordinates.json SHA256 b8671b1abbca5c7a406147dfad2359f1c1666887fab0b4eb67e489fc80abbe99
% Input: analyses/amp_transfer/inputs/redesign/refit-306-floor-endpoint.json SHA256 6d26e1e15c89217c7e29ce1d0772aa44285158071fdf4cc1bd0f6e7fb8fb2107
% Input: analyses/amp_transfer/notebook.qmd SHA256 a46134ecc164b9304778c6d809851a64a128dd91f932e6b898da3ce04bc9fd62
% Input: analyses/amp_transfer/results/kundu/summary.json SHA256 814f2ce5a28047c470bbcab59aadf13984e7ae034b55a4239b8e5b3058d0b55e
% Input: analyses/amp_transfer/results/refit-306/kundu-pressures.csv SHA256 561d08cf41e18d6ac277ee5e2dd207904a8885274d8d45739a0ec3f56d851cdb
% Input: analyses/amp_transfer/results/refit-306/refit-306-floor-fit.json SHA256 d7ee97cb1550d2a6969de406f09048dd4326f025a79523b7e205a5c65439233c
% Input: analyses/amp_transfer/results/refit-306/summary.json SHA256 3912f54eb7821b08e05f902d3001493b269355bea383439cada4e33fc4b4128f
% Input: analyses/amp_transfer/results/refit-306/uncertainty/prediction-bands.csv SHA256 d65f885c32fc9b391305b7e0ac76ec9f0112435455d027391c66dcd1ba7a9ff0
% Input: analyses/amp_transfer/results/refit-306/uncertainty/ranking.csv SHA256 ce80d211f43c53c9f654ec85ebdd5a91ac779a827c5e992415954cacffcab2a6
% Input: analyses/dea_transfer/inputs/runs/B-2-coordinates.json SHA256 7402f93dd3334eae251b38df951d2fdc310a730bf0045871bb83482524e75beb
% Input: analyses/dea_transfer/inputs/runs/B-2-endpoint-declarations.toml SHA256 a0e6cd9a5a7796d892130fc71f756480465898fa28cb0fc8cd371e9f4b730b72
% Input: analyses/dea_transfer/inputs/runs/B-2-endpoint-document.json SHA256 91c2384ef9f7b2a035962aafb72d97435a18a94595c5f16c64c9e2ac18785cdd
% Input: analyses/dea_transfer/notebook.qmd SHA256 c08a11c9e055d810c4b04c5c421a3ac89cf9edea895c7eb2ce3225e22368373d
% Input: analyses/dea_transfer/results/figure-values.csv SHA256 58cbb7b6e6f825b0f4fba1ef19f5393ad469e8aa42a58ece391f3b3d7af4bf2c
% Input: analyses/dea_transfer/results/fits.json SHA256 b936638b0b1f83802809bea15c003c1bbcc229b23824ce61d48ac4123b80af48
% Input: analyses/dea_transfer/results/score.json SHA256 12aa281fa03ed5d34cabbe3a09174b49da105199b15f1b16087ffa33018fd64e
% Input: analyses/dea_transfer/results/uncertainty-307/prediction-bands.csv SHA256 b59d9515e2d86b3c76eea92e97343eb5b2e28512bf4e8dd9f69e359de593f9ed
% Input: analyses/dea_transfer/results/uncertainty-307/ranking.csv SHA256 14ec3184971afe521ae4be0f1d636acd3a47067df1a01feb6068d75cf63cf155
% Input: analyses/mdea_mea_workflow/notebook.qmd SHA256 55831c2a23404765f14fbcbeb2c83eb64b7245a10730ad9ad9329b85aa1c2533
% Input: analyses/mdea_mea_workflow/results/refit-298/B1-endpoint/endpoint.json SHA256 2a5341be176829d087133a1467828f762a655f92a586eb1062a5719a58142d9d
% Input: analyses/mdea_mea_workflow/results/refit-298/B1-endpoint/fitted-replay.json SHA256 6690d0b7baf4bb3f08a9437743857d965d7a4a4fd1529d4e6f8e5d234f4059e2
% Input: analyses/mdea_mea_workflow/results/refit-298/attempts.json SHA256 40e17071831a699fe380d4ac4b1b286d5d283ad434a0e704ac2803d7d22c954f
% Input: analyses/mdea_mea_workflow/results/refit-298/pressure-comparison.csv SHA256 d2cdaf4a3bc1a2b7b9b957585a6dbaec2eff15ca56c009b862dac63a2912c22c
% Input: analyses/mea_parameter_bundle/results/selected-current-best-parameters.json SHA256 f8f64820b161377a0d97d1e4ad4f6cea764cd0558dc63231d58e2096988eef97
% Input: analyses/pz_transfer/notebook.qmd SHA256 a71c49071169b4a2795c538b22d2e59c58870489e0131cecf480d5fc98622858
% Input: analyses/pz_transfer/results/holdout-A12.csv SHA256 89f6aabc166a2663ddf7be4cfdc20821cc3ef9ef46070535644f5cd2abb66b15
% Input: analyses/pz_transfer/results/reactive-A12-record.json SHA256 3eefe76e4599050e524b9a1a8354121caae0ca30de38613359cf7e1f9edc1329
% Input: analyses/pz_transfer/results/reactive-A12-residuals.csv SHA256 3c8855bba6dfaa8ed8c8e64f7c4a160cafb77fb1d85b6858e8d5520e1469e4e4
% Input: analyses/pz_transfer/results/reactive-A12-start2-endpoint.json SHA256 c24ff09badc642d5dd312f5dc03248f2e5fa4c93e3b1d91fa69fe42120297b9e
% Input: analyses/pz_transfer/results/speciation-reactions.json SHA256 64a51633b7d56ce75a1246fe51eb6a39a1f0db1abc7d75b16b2fe13df9c125be
% Input: analyses/pz_transfer/results/uncertainty-A12/ranking.csv SHA256 f266707b75ee505bb8a5226c79e6f30037696720f859ca7d239972e6f58c71fe
% Input: analyses/solvent_specification_mdea/B1-declarations.toml SHA256 7b5fd083d83731c74624707f6e9f92bad05265c9ed19d105b25cf99dd897a042
% Input: analyses/solvent_specification_mdea/inputs/B1-coordinates.json SHA256 0bca9af202313b2bd8325ccaa37a72f32a1b326099c957c5ebf35f630506c8be
% Input: analyses/solvent_specification_mdea/notebook.qmd SHA256 227db268f1475ce3e11e4873c68e3337bf217918fa0ce0bdc11deeb25f57ff88
% Input: analyses/solvent_specification_mdea/results/adoption-302/B1/pressure.json SHA256 b81bfeb21d211d64211a2bda07b5e07af2a897f828a750f01aac6142340a7820
% Input: analyses/solvent_specification_mdea/results/retained-statistics-320.json SHA256 64393229910bf06ba614d1345e1fd2d2f675e2aba1a3a9a7eaedcd33ffd07e81
% Input: analyses/solvent_specification_mdea/results/uncertainty-307/prediction-bands.csv SHA256 96fc9352b6f30dae38805b799344404528d6e49b3b45eaa42002daa8523ac0cd
% Input: analyses/solvent_specification_mdea/results/uncertainty-307/predictions.json SHA256 8fd82d94b500a2b8985ea9979e4ceb653fbfb781dc66a3d047ee68bdec0f46ba
% Input: analyses/solvent_specification_mdea/results/uncertainty-307/ranking.csv SHA256 16b4ed6165c0cc7bad4d938567c8f8e8d88bdddb3de4eb8b53c43b372f7c8e7e
% Input: docs/scientific/latex/references.bib SHA256 7dc97086df08bf802d25d2a588158a9ad3881648729bb9a5da84741ef50acd9c
% Input: docs/scientific/latex/supplement.tex SHA256 b90143f6a0711b9a7437b4845c8a5db430f8b301ff21d8f402821b105da3d5f8
\begingroup\small\setlength{\tabcolsep}{3pt}
\begin{longtable}{@{}>{\raggedright\arraybackslash}p{0.32\linewidth}>{\raggedright\arraybackslash}p{0.36\linewidth}>{\raggedright\arraybackslash}p{0.25\linewidth}@{}}
\caption{Transferred quantities and their origins. All shared quantities remain fixed here; amine-specific interactions with shared species may be fitted.}\label{tab:p2-origins}\\
\toprule Quantity & Origin in Part I & Implication here \\ \midrule\endfirsthead
\toprule Quantity & Origin in Part I & Implication here \\ \midrule\endhead
\bottomrule\endfoot
Water segment, energy and association coefficients & Independent fit to IAPWS-95 & Fixed; not fitted to any of these amines \\
\ce{CO2} segment and energy coefficients & Gross and Sadowski, Table 2 \cite{Gross2001} & Literature inputs; fixed \\
\ce{CO2}--water dispersion correction and slope & Independent fit to Kiepe pressures & Fixed \\
Inorganic-ion segment sizes and energies & Uyan et al., Table 4 \cite{Uyan2015} & Fixed; hydrogen-ion values used for hydronium \\
Inorganic Born sizes and remaining water--ion terms & Assumed inputs in Part I & Fixed; not independently measured in MEA \\
Bicarbonate--water interaction & MEA joint fit, upper bound $+1$ & Fixed at bound; conditional SE excludes this direction \\
Completed R2 reaction level and enthalpy correction & MEA joint fit; interior with conditional errors & Fixed; source-law correction, not a new measurement \\
R1 and R3 water/carbonate laws & Literature laws as mapped in Part I & Fixed with source-reference conversion \\
Solvation and dielectric descriptions & Literature coefficients and assumptions in Part I & Water/CO2 inputs fixed; amine coefficients specified separately \\
Reactive coordinates (MDEA / AMP / PZ / DEA) & Amine-specific fits: 9 / 9 / 12 / 4 & Not totals of earlier pure and binary fitting \\
\end{longtable}\endgroup


% P:p2-model-02
Each liquid contains the shared species and its molecular amine and protonated forms.
AMP and DEA also contain carbamate; PZ contains singly and doubly protonated PZ, mono- and dicarbamate, and the neutral zwitterion HPZCOO.
The final reactive fits have nine free coordinates for MDEA, nine for AMP, twelve for PZ and four for DEA, excluding earlier pure-component, binary and dielectric fits.
Supplementary Section~S1 gives species parameters, site descriptions, diameter ties and temperature laws.
The fitted ion diameters and reaction-law corrections are effective model parameters, not independently measured molecular sizes or equilibrium constants.

% P:p2-model-03
Charged species remain liquid-only.
The vapor contains water, \COtwo and molecular amine, and for PZ also HPZCOO.
Allowing this neutral zwitterion in the vapor resolves otherwise undecided phase searches; at the earlier fixed PZ fit, the relative change in fitted partial pressures is at most $1.2\times10^{-8}$.
This numerical change is negligible for those pressures, but the vapor composition space matters to the phase search.
PZ has no solid-phase description, and its pure-liquid density is unvalidated.

\subsection{Reaction conventions and measured temperature ranges}\label{sec:p2-laws}
% P:p2-model-04
The pure-liquid solvent and infinite-dilution solute conventions follow Part I.
The amine laws and shared R2 retain their source molality standard state; R1 and R3 retain the source mole-fraction basis.
Explicit water activity and the same hydronium reference are used, with the conversions described in Part I.
For a reaction written with stoichiometric coefficients $\nu_i$ positive for products, equilibrium requires
\begin{equation}
 \sum_i\nu_i\mu_i=0,\qquad \ln K(T)=a+b/T+c\ln(T/T_0)+dT+\cdots .
 \label{eq:p2-law}
\end{equation}
Source constants are mapped to this reference once; mole-fraction and molality constants are not interchanged without their standard-state conversion.
Declared calculation intervals and measured source intervals answer different questions, as shown in Table~\ref{tab:p2-laws}.

% Generator: docs/scientific/latex/scripts/part2_displays.py; no model evaluation.
% Input: analyses/amp_transfer/inputs/kundu/current-declarations.toml SHA256 3085081adcf6b1d5e797a49ae6b96037cc1b5bc5191f45cff699d8546d774c55
% Input: analyses/amp_transfer/inputs/redesign/refit-306-floor-coordinates.json SHA256 b8671b1abbca5c7a406147dfad2359f1c1666887fab0b4eb67e489fc80abbe99
% Input: analyses/amp_transfer/inputs/redesign/refit-306-floor-endpoint.json SHA256 6d26e1e15c89217c7e29ce1d0772aa44285158071fdf4cc1bd0f6e7fb8fb2107
% Input: analyses/amp_transfer/notebook.qmd SHA256 a46134ecc164b9304778c6d809851a64a128dd91f932e6b898da3ce04bc9fd62
% Input: analyses/amp_transfer/results/kundu/summary.json SHA256 814f2ce5a28047c470bbcab59aadf13984e7ae034b55a4239b8e5b3058d0b55e
% Input: analyses/amp_transfer/results/refit-306/kundu-pressures.csv SHA256 561d08cf41e18d6ac277ee5e2dd207904a8885274d8d45739a0ec3f56d851cdb
% Input: analyses/amp_transfer/results/refit-306/refit-306-floor-fit.json SHA256 d7ee97cb1550d2a6969de406f09048dd4326f025a79523b7e205a5c65439233c
% Input: analyses/amp_transfer/results/refit-306/summary.json SHA256 3912f54eb7821b08e05f902d3001493b269355bea383439cada4e33fc4b4128f
% Input: analyses/amp_transfer/results/refit-306/uncertainty/prediction-bands.csv SHA256 d65f885c32fc9b391305b7e0ac76ec9f0112435455d027391c66dcd1ba7a9ff0
% Input: analyses/amp_transfer/results/refit-306/uncertainty/ranking.csv SHA256 ce80d211f43c53c9f654ec85ebdd5a91ac779a827c5e992415954cacffcab2a6
% Input: analyses/dea_transfer/inputs/runs/B-2-coordinates.json SHA256 7402f93dd3334eae251b38df951d2fdc310a730bf0045871bb83482524e75beb
% Input: analyses/dea_transfer/inputs/runs/B-2-endpoint-declarations.toml SHA256 a0e6cd9a5a7796d892130fc71f756480465898fa28cb0fc8cd371e9f4b730b72
% Input: analyses/dea_transfer/inputs/runs/B-2-endpoint-document.json SHA256 91c2384ef9f7b2a035962aafb72d97435a18a94595c5f16c64c9e2ac18785cdd
% Input: analyses/dea_transfer/notebook.qmd SHA256 c08a11c9e055d810c4b04c5c421a3ac89cf9edea895c7eb2ce3225e22368373d
% Input: analyses/dea_transfer/results/figure-values.csv SHA256 58cbb7b6e6f825b0f4fba1ef19f5393ad469e8aa42a58ece391f3b3d7af4bf2c
% Input: analyses/dea_transfer/results/fits.json SHA256 b936638b0b1f83802809bea15c003c1bbcc229b23824ce61d48ac4123b80af48
% Input: analyses/dea_transfer/results/score.json SHA256 12aa281fa03ed5d34cabbe3a09174b49da105199b15f1b16087ffa33018fd64e
% Input: analyses/dea_transfer/results/uncertainty-307/prediction-bands.csv SHA256 b59d9515e2d86b3c76eea92e97343eb5b2e28512bf4e8dd9f69e359de593f9ed
% Input: analyses/dea_transfer/results/uncertainty-307/ranking.csv SHA256 14ec3184971afe521ae4be0f1d636acd3a47067df1a01feb6068d75cf63cf155
% Input: analyses/mdea_mea_workflow/notebook.qmd SHA256 55831c2a23404765f14fbcbeb2c83eb64b7245a10730ad9ad9329b85aa1c2533
% Input: analyses/mdea_mea_workflow/results/refit-298/B1-endpoint/endpoint.json SHA256 2a5341be176829d087133a1467828f762a655f92a586eb1062a5719a58142d9d
% Input: analyses/mdea_mea_workflow/results/refit-298/B1-endpoint/fitted-replay.json SHA256 6690d0b7baf4bb3f08a9437743857d965d7a4a4fd1529d4e6f8e5d234f4059e2
% Input: analyses/mdea_mea_workflow/results/refit-298/attempts.json SHA256 40e17071831a699fe380d4ac4b1b286d5d283ad434a0e704ac2803d7d22c954f
% Input: analyses/mdea_mea_workflow/results/refit-298/pressure-comparison.csv SHA256 d2cdaf4a3bc1a2b7b9b957585a6dbaec2eff15ca56c009b862dac63a2912c22c
% Input: analyses/mea_parameter_bundle/results/selected-current-best-parameters.json SHA256 f8f64820b161377a0d97d1e4ad4f6cea764cd0558dc63231d58e2096988eef97
% Input: analyses/pz_transfer/notebook.qmd SHA256 a71c49071169b4a2795c538b22d2e59c58870489e0131cecf480d5fc98622858
% Input: analyses/pz_transfer/results/holdout-A12.csv SHA256 89f6aabc166a2663ddf7be4cfdc20821cc3ef9ef46070535644f5cd2abb66b15
% Input: analyses/pz_transfer/results/reactive-A12-record.json SHA256 3eefe76e4599050e524b9a1a8354121caae0ca30de38613359cf7e1f9edc1329
% Input: analyses/pz_transfer/results/reactive-A12-residuals.csv SHA256 3c8855bba6dfaa8ed8c8e64f7c4a160cafb77fb1d85b6858e8d5520e1469e4e4
% Input: analyses/pz_transfer/results/reactive-A12-start2-endpoint.json SHA256 c24ff09badc642d5dd312f5dc03248f2e5fa4c93e3b1d91fa69fe42120297b9e
% Input: analyses/pz_transfer/results/speciation-reactions.json SHA256 64a51633b7d56ce75a1246fe51eb6a39a1f0db1abc7d75b16b2fe13df9c125be
% Input: analyses/pz_transfer/results/uncertainty-A12/ranking.csv SHA256 f266707b75ee505bb8a5226c79e6f30037696720f859ca7d239972e6f58c71fe
% Input: analyses/solvent_specification_mdea/B1-declarations.toml SHA256 7b5fd083d83731c74624707f6e9f92bad05265c9ed19d105b25cf99dd897a042
% Input: analyses/solvent_specification_mdea/inputs/B1-coordinates.json SHA256 0bca9af202313b2bd8325ccaa37a72f32a1b326099c957c5ebf35f630506c8be
% Input: analyses/solvent_specification_mdea/notebook.qmd SHA256 227db268f1475ce3e11e4873c68e3337bf217918fa0ce0bdc11deeb25f57ff88
% Input: analyses/solvent_specification_mdea/results/adoption-302/B1/pressure.json SHA256 b81bfeb21d211d64211a2bda07b5e07af2a897f828a750f01aac6142340a7820
% Input: analyses/solvent_specification_mdea/results/retained-statistics-320.json SHA256 64393229910bf06ba614d1345e1fd2d2f675e2aba1a3a9a7eaedcd33ffd07e81
% Input: analyses/solvent_specification_mdea/results/uncertainty-307/prediction-bands.csv SHA256 96fc9352b6f30dae38805b799344404528d6e49b3b45eaa42002daa8523ac0cd
% Input: analyses/solvent_specification_mdea/results/uncertainty-307/predictions.json SHA256 8fd82d94b500a2b8985ea9979e4ceb653fbfb781dc66a3d047ee68bdec0f46ba
% Input: analyses/solvent_specification_mdea/results/uncertainty-307/ranking.csv SHA256 16b4ed6165c0cc7bad4d938567c8f8e8d88bdddb3de4eb8b53c43b372f7c8e7e
% Input: docs/scientific/latex/references.bib SHA256 9ac4ac1c25d01550aa5e9a9df066b7a49e96212687fb4c6f9c826ee81cea4252
% Input: docs/scientific/latex/supplement.tex SHA256 b90143f6a0711b9a7437b4845c8a5db430f8b301ff21d8f402821b105da3d5f8
\begingroup\small\setlength{\tabcolsep}{3pt}
\begin{longtable}{@{}>{\raggedright\arraybackslash}p{0.18\linewidth}>{\raggedright\arraybackslash}p{0.23\linewidth}>{\raggedright\arraybackslash}p{0.22\linewidth}>{\raggedright\arraybackslash}p{0.28\linewidth}@{}}
\caption{Reaction-law measurement intervals and use. A calculation inside a declared interval can still extrapolate the measurements. Rounded source limits are distinguished from actual tabulated temperatures.}\label{tab:p2-laws}\\
\toprule Law & Measured interval / K & Declared interval / K & Use and limitation \\ \midrule\endfirsthead
\toprule Law & Measured interval / K & Declared interval / K & Use and limitation \\ \midrule\endhead
\bottomrule\endfoot
MDEA dissociation & 278--368 & 278--395 & 393.15 K pressure and 388.71 K heat extrapolate measurements \\
AMP dissociation & to 353.22 & to 353.22 & Tong 373--393 K rows not calculated (outside declared interval) \\
AMP hydrolysis & source slope: 313--353 & includes 303--353 & Slope extrapolated below 313 K; level and slope fitted \\
PZ protonation & approximately 293--353 & includes 283.15--393.15 & 283.15 and 393.15 K extrapolate \\
PZ carbamate & 283.15--333.15 (rounded 283--333) & includes 353.15 and 393.15 & 353.15 and 393.15 K extrapolate \\
DEA protonation & to approximately 323 & to 393.15 & 23/30 fitted Lawson rows above measured range \\
Shared R2 & not established here & 293.15--393.15 & Seven Lawson rows at 394.26 K not calculated (outside declared interval); PZ 283.15 K extrapolates below lower limit \\
\end{longtable}\endgroup


% P:p2-model-05
The MDEA dissociation coefficients are those of Kamps and Maurer, measured over 278--368 K, although the computational declaration extends to 395 K.
The fixed AMP dissociation law ends at 353.22 K.
Its carbamate-hydrolysis temperature dependence is fitted: the completed $b$ is $-1944.8$ K, compared with the starting source slope $-2546.2$ K, while the enthalpy correction reaches its $-5000$ J/mol limit.
This fitted hydrolysis slope is distinct from the fixed dissociation law and does not extend that law's admissible interval.
PZ carbamate source measurements span 283.15--333.15 K, and its protonation measurements approximately 293--353 K.
DEA uses the fixed Bower protonation law above the approximately 323 K upper measurement limit, including 23 of its 30 fitted pressures.
A fitted offset does not turn an extrapolated temperature dependence into a measured one.
